% !TeX root = ../../Book2.tex
\section{乘子算子}
\label{b2:sec:2201}

平移不变算子在 Fourier 一侧表现为逐点乘法。本章讨论怎样从符号的局部光滑性得到整个 \(L^p\) 范数界，并区别阶数为零的有界乘子、改变可积指数的分数积分和具有导数损失的微分算子。

\subsection{Fourier 乘子}
\label{b2:subsec:220101}

\begin{definition}\label{b2:def:22-multiplier}
对可测函数 \(m\)，先在使乘积有意义的测试函数上定义
\[
 T_m f=\mathcal F^{-1}(m\widehat f).
\]
若该算子延拓为 \(L^p(\mathbb R^d)\) 上的有界算子，称 \(m\) 为 \(L^p\) Fourier 乘子，并记其算子范数为 \(\|m\|_{M_p}\)。
\end{definition}
乘子在零测集上的值不影响算子。对于一般有多项式增长的光滑 \(m\)，乘法可在 Schwartz 空间与温和分布之间解释；对于仅可测的有界 \(m\)，最稳妥的起点是 \(L^2\)。
若 \(m=\widehat k\)、\(k\in L^1\)，则 \(T_mf=k*f\)，Young 不等式给出所有 \(1\le p\le\infty\) 上的有界性。

\ProseParagraphBreak
乘法和卷积的对应不意味着每个乘子都有普通可积核。恒等乘子的核为 \(\delta_0\)，Hilbert 乘子的核为主值分布。推导核估计时应先局部化频率，再在原点外重构普通函数；原点上的分布项要由已有的 \(L^2\) 定义处理。

\subsection{\texorpdfstring{\(L^2\)}{L2} 理论}
\label{b2:subsec:220102}

\begin{proposition}\label{b2:prop:22-l2}
可测 \(m\) 是 \(L^2\) 乘子当且仅当 \(m\in L^\infty\)，且
\(\|m\|_{M_2}=\|m\|_\infty\)。
\end{proposition}
\begin{proof}
Plancherel 给出充分性及范数上界。若算子范数为 \(C\)，在任意有限测度集合 \(E\) 上选 \(\widehat f=\mathbf1_E\)，得
\(\int_E|m|^2\le C^2|E|\)。若 \(|m|>C+\epsilon\) 的集合有正测度，取其中有限测度的正测度子集就矛盾。因此 \(|m|\le C\) 几乎处处，得到必要性和等号。
\end{proof}
对 \(|m|=1\) 的符号，算子在 \(L^2\) 上等距；这并不提供其他 \(p\) 上的范数控制。即使每个频率只发生相位旋转，不同频率叠加后的空间集中仍可能发生变化。

\ProseParagraphBreak
若 \(m_1,m_2\in M_p\)，则 \(m_1m_2\in M_p\)，且范数不超过两范数之积。先在 \(L^2\cap L^p\) 上验证 \(T_{m_1m_2}=T_{m_1}T_{m_2}\)，再用稠密性延拓即可。相似地，乘子空间对一致算子范数有界的点态极限封闭：在 \(L^2\cap L^p\) 上先取 \(L^2\) 极限及几乎处处子列，再用 Fatou，最后延拓。这是处理光滑截断极限的常用办法。

\subsection{\texorpdfstring{\(L^p\)}{Lp} 乘子问题}
\label{b2:subsec:220103}

若 \(m\in M_p\)、\(1<p<\infty\)，其 Banach 伴随在 \(L^{p'}\) 上的符号为 \(\overline m\)，范数相同。若 \(m\) 实值，就得到对偶指数上的相同范数。与 \(L^2\) 界插值，可控制 \(p\) 和 \(2\) 之间的全部指数，但不能越过尚未建立的端点。

\ProseParagraphBreak
缩放不改变乘子范数。若 \(D_rf(x)=f(rx)\)，则
\(T_{m(r\,\cdot)}=D_{1/r}T_mD_r\)，而两个缩放范数相消。
正交变换与频率平移也保持范数：前者对应空间坐标旋转，后者对应调制。符号乘上 \(e^{-2\pi\mathrm i y\cdot\xi}\) 则对应输出平移。这些等距对称使乘子条件应具有尺度不变性；要求所有导数在整个频率空间绝对有界，并不是适合齐次问题的条件。

\ProseParagraphBreak
一个不光滑却有界的例子是一维区间特征函数。半轴投影可用恒等算子与 Hilbert 变换表示，两个频率平移后的半轴投影之差就是区间投影，因此在 \(1<p<\infty\) 上一致有界。由此不能推断高维球的特征函数同样有界。符号边界的几何维数与曲率会参与问题，\chapref{b2:ch:24}将专门说明这一障碍。

% 扩写定位：b2:subsec:220103
Fourier 乘子并不是凭空选出的一类算子；在 \(L^2\) 上，它们正好是与所有平移交换的有界线性算子。
\begin{proposition}\label{b2:prop:22-translation-multiplier}
若 \(T:L^2(\mathbb R^d)\to L^2(\mathbb R^d)\) 有界，且 \(T\tau_y=\tau_yT\) 对全部 \(y\) 成立，其中 \(\tau_yf(x)=f(x-y)\)，则存在唯一的 \(m\in L^\infty\)，使 \(T=T_m\)，并且 \(\|T\|_{2\to2}=\|m\|_\infty\)。
\end{proposition}
\begin{proof}
令 \(A=\mathcal FT\mathcal F^{-1}\)。平移变成乘法 \(e^{-2\pi\mathrm iy\cdot\xi}\)，故 \(A\) 与所有这些乘法算子交换。对任意 \(\varphi\in C_c^\infty\)，Fourier 反演把 \(\varphi\) 写成这些指数乘法的 \(L^1\) 加权积分；其在 \(L^2\) 上的强积分与 \(A\) 交换。于是 \(A\) 与乘法 \(M_\varphi\) 交换。

由光滑函数的有界几乎处逼近，\(A\) 也与任意可测集合的投影 \(M_{\mathbf1_E}\) 交换。可以先对有界集合用 Lebesgue 正则性与卷积逼近得到 \(0\leq\varphi_n\leq1\)、\(\varphi_n\to\mathbf1_E\) 几乎处，支配收敛给出强算子收敛；再以球截断推广到一般集合。

取递增的有限测度立方体 \(Q_n\) 覆盖全空间。在 \(Q_n\) 上定义 \(m_n=A\mathbf1_{Q_n}\)。若 \(n<k\)，交换投影给出
\(\mathbf1_{Q_n}m_k=A(\mathbf1_{Q_n}\mathbf1_{Q_k})=m_n\)，所以拼成可测函数 \(m\)。对 \(E\subset Q_n\)，有
\[
 A\mathbf1_E=\mathbf1_E m,\qquad
 \int_E|m|^2\leq\|A\|^2|E|.
\]
若 \(\{|m|>\|A\|+\epsilon\}\cap Q_n\) 有正测度，将其代入即矛盾。因此 \(\|m\|_\infty\leq\|A\|\)。有限支集简单函数的稠密性说明 \(A=M_m\)，反向范数与唯一性由本节的 \(L^2\) 乘子结论得到。
\end{proof}
该结论只建立在 \(L^2\) 上。若算子只在 \(L^p\) 给定，还不能通过一次 Fourier 变换就把它视为 \(L^2\) 有界乘子；需先说明它在共同稠密类上的作用及其额外的平方界。另一方面，具有不连续符号的平移不变算子在 \(L^2\) 上完全没有困难；真正的乘子问题是同一个算子能否保持其他 \(L^p\)。
